\documentclass[12pt]{article}
\usepackage{amsfonts,amsthm,amsmath,amssymb,mathtools}
\usepackage{mathrsfs}
\usepackage{enumitem}
\usepackage{tikz-cd}
\usepackage{geometry}
\usepackage{mlmodern}
\usepackage{xcolor}
\usepackage{pgfplots}

\geometry{letterpaper, portrait, margin=1in}
\pgfplotsset{compat=1.18}

\setcounter{section}{0}

\renewcommand{\thesection}{\S\arabic{section}}
\renewcommand{\thesubsection}{\arabic{section}}

\newtheorem{thm}{Theorem}
\newenvironment{theorem}[1]{
	\renewcommand{\thethm}{#1}
	\begin{thm}
		}{\end{thm}}

\newtheorem{lma}{Lemma}
\newenvironment{lemma}[1]{
	\renewcommand{\thelma}{#1}
	\begin{lma}
		}{\end{lma}}

\theoremstyle{definition}
\newtheorem{ex}{}
\newenvironment{exercise}[1]{
	\renewcommand{\theex}{#1}
	\begin{ex}
		}{\end{ex}}

\title{Homework Assignment 4}
\author{Connor Mooneyhan}
\date{September 19, 2025}

\begin{document}

\maketitle

\begin{ex}
	(2B.23)
	Suppose $f: \mathbb{R} \to \mathbb{R}$ is a strictly increasing function.
	Prove that the inverse function $f^{-1} : f(\mathbb{R}) \to \mathbb{R}$ is a continuous function.
	\newline[\textit{Note that this exercise does not have as a hypothesis that $f$ is continuous}]
\end{ex}
\begin{proof}
	Let $U$ be open in $\mathbb{R}$.
	We want to show that $(f^{-1})^{-1}(U) = f(U)$ is open in $f(\mathbb{R})$.
	Let $x \in f(U)$.
	Then let $(a, b) \subset \mathbb{R}$ such that $f^{-1}(x) \in (a, b) \subset U$.
	Since $f$ is strictly increasing and $a < f^{-1}(x) < b$, $f(a) < x < f(b)$.
	In particular, $x \in (f(a), f(b)) \subset f(U)$, which is open in the subspace/order topology on $f(\mathbb{R})$.
	Thus every point of $f(U)$ has a neighborhood in $f(\mathbb{R})$ also contained in $U$, so $f(U)$ is open.
\end{proof}

\begin{ex}
	(2B.24)
	Suppose $B \subset \mathbb{R}$ is a Borel set and $f: B \to \mathbb{R}$ is an increasing function.
	Prove that $f(B)$ is a Borel set.
\end{ex}
\begin{proof}
  This one is escaping me.
  Given that it comes right after the previous problem and has a similar setup, I tried especially hard to use it.
  One option I thought was promising was to turn $f$ into a strictly increasing function by considering the function $g : B/\sim \,\to \mathbb{R}$ where $B/\sim$ is the collection of fibers over $g$.
  Indeed, $g$ is strictly increasing, but unfortunately it is no longer a function from $\mathbb{R}$ to $\mathbb{R}$ since the elements of the domain are subsets of $\mathbb{R}$.
  I tried to still recover the core idea, considering open sets in the quotient topology on $B/\sim$ and trying to figure out what they could tell me about how open sets map in $f$, but I haven't been able to figure it out.

  Another idea included considering when $B$ is an open set, but nothing ocurred to me that made that case any easier.
  If it had worked, I might have been able to make a similar argument about closed sets, then said something about countable unions (resp. intersections) of closed (resp. open) sets.

  I think the difficulty lies in proving things about Borel sets directly.
  If it were a \textit{collection} of potentially Borel sets, then I might be able to do something with showing that collection was a $\sigma$-algebra that contained open sets, but I don't see any room for that sort of argument here.
\end{proof}

\begin{ex}
	(2C.5)
	Suppose $(X, S, \mu)$ is a mesaure space such that $\mu(X) < \infty$.
	Prove that if $\mathcal{A}$ is a set of disjoint sets in $S$ such that $\mu(A) > 0$ for every $A \in \mathcal{A}$, then $\mathcal{A}$ is a countable set.
\end{ex}
\begin{proof}
	Let $\mu(X) = k$, and define $\mathcal{A}_n = \left\lbrace A \in \mathcal{A} : \mu(A) > \frac{1}{n} \right\rbrace$ for all $n \in \mathbb{Z^+}$.
	Then \begin{equation*}
		k \geq \mu \left( \bigcup_{A \in A_n} A \right) \geq \mu \left( \bigcup_{i = 1}^N A_i \right) = \sum_{i = 1}^N \mu(A) \geq \frac{N}{n}.
	\end{equation*}
	where the first and second inequalities come from monotonicity and the equality comes from finite additivity.
	Note that the right hand side of the first inequality just represents any finite sum of elements of $\mathcal{A}_n$.
	If $\mathcal{A}_n$ is empty, then we can take the right hand side of each inequality to be 0.
	Thus, we cannot be allowed to take $N > kn$, so $\mathcal{A}_n$ must have at most $kn$ elements, and in particular, finitely many.
	Then, since $\mathcal{A} = \bigcup_{n \in \mathbb{Z}^+} \mathcal{A}_n$ is a countable union of finite sets, $\mathcal{A}$ is countable.
\end{proof}

\begin{ex}
	(2C.9)
	Suppose $\mu$ and $\nu$ are measures on a measurable space $(X, S)$.
	Prove that $\mu + \nu$ is a measure on $(X, S)$.
		[Here $\mu + \nu$ is the usual sum of two functions: if $E \in S$, then $(\mu + \nu)(E) = \mu(E) + \nu(E)$.]
\end{ex}
\begin{proof}
	First, $(\mu + \nu)(\varnothing) = \mu(\varnothing) + \nu(\varnothing) = 0 + 0 = 0$.
	Now let $\mathcal{A} \subset S$ be a countable collection of measurable sets.
	Then \begin{align*}
		(\mu + \nu)\left( \bigcup_{A \in \mathcal{A}} A \right) & = \mu\left( \bigcup_{A \in \mathcal{A}} A \right) + \nu\left( \bigcup_{A \in \mathcal{A}} A \right) \\
		                                                        & = \sum_{A \in \mathcal{A}}\mu(A) + \sum_{A \in \mathcal{A}}\nu(A)                                   \\
		                                                        & = \sum_{A \in \mathcal{A}}\mu(A) + \nu(A)                                                           \\
		                                                        & = \sum_{A \in \mathcal{A}}(\mu + \nu)(A)
	\end{align*}
	as desired.
\end{proof}

\begin{ex}
	(2C.11)
	Suppose $(X, S, \mu)$ is a measure space and $C, D, E \in S$ are such that \begin{equation*}
		\mu(C \cap D) < \infty, \mu(C \cap E) < \infty, \text{ and } \mu(D \cap E) < \infty.
	\end{equation*}
	Find and prove a formula for $\mu(C \cup D \cup E)$ in terms of $\mu(C)$, $\mu(D)$, $\mu(E)$, $\mu(C \cap D)$, $\mu(C \cap E)$, $\mu(D \cap E)$, and $\mu(C \cap D \cap E)$.
\end{ex}
\begin{proof}
	Our proof will use Theorem 2.61 (the formula for the measure of the union of two sets whose union is finite) extensively.
	Two more facts to note are that \begin{equation*}
		\mu((C \cup D) \cap E) = \mu((C \cap E) \cup (D \cap E)) \leq \mu(C \cap E) + \mu(D \cap E) < \infty
	\end{equation*}
  by subadditivity and \begin{equation*}
    \mu((C \cap E) \cap (D \cap E)) = \mu(C \cap D \cap E) \leq \mu(C \cap D) < \infty
  \end{equation*}
  by monotonicity.
	Now, we see that \begin{align*}
		\mu(C \cup D \cup E) & = \mu((C \cup D) \cup E)                                                                                       \\
		                     & = \mu(C \cup D) + \mu(E) - \mu((C \cup D) \cap E)                                                              \\
		                     & = \mu(C) + \mu(D) - \mu(C \cap D) + \mu(E) - \mu((C \cap E) \cup (D \cap E))                                   \\
		                     & = \mu(C) + \mu(D) - \mu(C \cap D) + \mu(E) - (\mu(C \cap E) + \mu(D \cap E) - \mu((C \cap E) \cap (D \cap E))) \\
                         & = \mu(C) + \mu(D) + \mu(E) - \mu(C \cap D) - \mu(C \cap E) - \mu(D \cap E) + \mu(C \cap D \cap E).
	\end{align*}
\end{proof}

\begin{ex}
	(2D.5)
	Prove that if $A \subset \mathbb{R}$ is Lebesgue measurable, then there exists an increasing sequence $F_1 \subset F_2 \subset \cdots$ of closed sets contained in $A$ such that \begin{equation*}
		\left| A \setminus \bigcup_{k=1}^\infty F_k \right| = 0.
	\end{equation*}
\end{ex}
\begin{proof}
  By Theorem 2.71, we can choose closed sets $G_1, G_2, \dots$ contained in $A$ such that $\left| A \setminus \bigcup_{k=1}^\infty G_k \right| = 0$.
  We want the sequence to be increasing, so define $F_n = \bigcup_{k=1}^n G_k$.
  Then $F_n \subset F_{n+1}$ for all $n$, and each $F_n$ is closed since it is a finite union of closed sets.
  Since $\bigcup_{k=1}^\infty F_k = \bigcup_{k=1}^\infty G_k$, we have \begin{equation*}
    \left| A \setminus \bigcup_{k=1}^\infty F_k \right| = \left| A \setminus \bigcup_{k=1}^\infty G_k \right| = 0.
  \end{equation*}
\end{proof}

\section*{}

\textbf{What topics or problems would you like to review before Written Test \#1?}\\
Riemann integration!
\\\\

\end{document}
